% =====================================================================
% BMED365: Medical Image Analysis
% Beamer presentation - Topic List B01-B06
% =====================================================================
\documentclass[aspectratio=169, 10pt]{beamer}

% =====================================================================
% PACKAGES
% =====================================================================
\usepackage[utf8]{inputenc}
\usepackage[T1]{fontenc}
\usepackage[english]{babel}
\usepackage{graphicx}
\usepackage{tikz}
\usetikzlibrary{shapes.geometric, arrows, positioning, calc}
\usepackage{booktabs}
\usepackage{amsmath}
\usepackage{fontawesome5}
\usepackage{hyperref}
\hypersetup{colorlinks=true, linkcolor=blue, urlcolor=blue}

% =====================================================================
% THEME AND COLORS
% =====================================================================
\usetheme{Madrid}
\usecolortheme{beaver}

% =====================================================================
% TITLE INFO
% =====================================================================
\title{Medical Image Analysis}
\subtitle{BMED365: Topic List B01--B06}
\author{BMED365}
\date{Spring 2026}

% =====================================================================
% DOCUMENT
% =====================================================================
\begin{document}

% Title page
\begin{frame}
    \titlepage
\end{frame}

% Table of Contents
\begin{frame}{Overview}
    \tableofcontents
\end{frame}

% =====================================================================
% SECTION: MRI Fundamentals
% =====================================================================
\section{MRI Fundamentals}

\begin{frame}{B01: Explain fundamental MRI principles at a high level}
    \textbf{Magnetic Resonance Imaging (MRI):}
    
    \vspace{0.3cm}
    \textbf{Physical principles (simplified):}
    \begin{enumerate}
        \item \textbf{Magnetization:} Strong magnetic field aligns hydrogen atoms (protons)
        \item \textbf{RF pulse:} Radiofrequency pulse ``tips'' the protons
        \item \textbf{Relaxation:} Protons return to equilibrium, emitting a signal
        \item \textbf{Detection:} Signal is measured and reconstructed into an image
    \end{enumerate}
    
    \vspace{0.3cm}
    \textbf{Why MRI is important in medicine:}
    \begin{columns}[T]
        \begin{column}{0.48\textwidth}
            \begin{itemize}
                \item No ionizing radiation
                \item Excellent soft tissue contrast
                \item Many different sequences/contrasts
            \end{itemize}
        \end{column}
        \begin{column}{0.48\textwidth}
            \begin{itemize}
                \item Can image function (fMRI)
                \item Gold standard for brain, joints, spine
                \item Multiparametric -- rich information
            \end{itemize}
        \end{column}
    \end{columns}
    
    \begin{block}{Relevance for AI}
        MRI generates large amounts of data - ideal for AI-based analysis, segmentation, and pattern recognition.
    \end{block}
\end{frame}

\begin{frame}{B02: Know about different MR sequences (T1, T1+Gd, T2, FLAIR)}
    \textbf{Common MR sequences and their contrast:}

    \vspace{0.2cm}
    \begin{center}
        \footnotesize
        \begin{tabular}{lp{4cm}p{4.5cm}}
            \toprule
            \textbf{Sequence} & \textbf{Characteristics} & \textbf{Clinical use} \\
            \midrule
            \textbf{T1-weighted} & CSF dark, fat bright, good anatomy & Anatomy \\
            \addlinespace
            \textbf{T1 + Gd} & Contrast shows blood-brain barrier breakdown & Active tumor, metastases \\
            \addlinespace
            \textbf{T2-weighted} & CSF bright, pathology often bright & Edema, inflammation, fluid \\
            \addlinespace
            \textbf{FLAIR} & T2-like, but CSF dark & Lesions near ventricles, MS plaques \\
            \addlinespace
            \textbf{DWI} & Diffusion of water molecules & Acute stroke (within minutes!) \\
            \bottomrule
        \end{tabular}
    \end{center}

    \vspace{0.2cm}
    \begin{alertblock}{\footnotesize For AI analysis}
        \footnotesize Multiparametric MRI {\tiny (multiple sequences)} provides complementary information that AI can leverage for better classification.
    \end{alertblock}
\end{frame}

% =====================================================================
% SECTION: Segmentation
% =====================================================================
\section{Segmentation}

\begin{frame}{B03: Describe segmentation in medical image analysis}
    \textbf{Segmentation = identify and delineate structures in images}
    
    \vspace{0.3cm}
    \textbf{Task:}
    \begin{itemize}
        \item Assign each pixel/voxel to a \textbf{class} (e.g. tumor, normal tissue, background)
        \item Produces a \textbf{segmentation map} (mask)
    \end{itemize}
    
    \vspace{0.3cm}
    \textbf{Types of segmentation:}
    \begin{columns}[T]
        \begin{column}{0.48\textwidth}
            \textbf{Semantic segmentation:}
            \begin{itemize}
                \item Classify each pixel
                \item All tumors = one class
            \end{itemize}
        \end{column}
        \begin{column}{0.48\textwidth}
            \textbf{Instance segmentation:}
            \begin{itemize}
                \item Distinguishes individual objects
                \item Tumor 1, Tumor 2, ...
            \end{itemize}
        \end{column}
    \end{columns}
    
    \vspace{0.3cm}
    \textbf{AI methods for segmentation:}
    \begin{itemize}
        \item \textbf{U-Net:} Classic encoder-decoder architecture (medical standard)
        \item \textbf{\href{https://github.com/MIC-DKFZ/nnUNet}{nnU-Net}:} Self-configuring, state-of-the-art
        \item \textbf{Segment Anything (SAM):} General-purpose segmentation
    \end{itemize}
\end{frame}

\begin{frame}{B04: Know about the BraTS challenge for brain tumor segmentation}
    \textbf{\href{https://www.synapse.org/brats}{BraTS} = Brain Tumor Segmentation Challenge}

    \vspace{0.2cm}
    \begin{columns}[T]
        \begin{column}{0.48\textwidth}
            \textbf{What is BraTS?}
            \begin{itemize}
                \item Annual international competition since 2012
                \item Benchmark for AI segmentation of brain tumors
                \item Multiparametric MRI (T1, T1+Gd, T2, FLAIR)
            \end{itemize}

            \vspace{0.2cm}
            \textbf{Segmentation task:}
            \begin{itemize}
                \item \textbf{Whole tumor (WT):} All tumor tissue
                \item \textbf{Tumor core (TC):} Active + necrotic
                \item \textbf{Enhancing tumor (ET):} Contrast-enhancing
            \end{itemize}
        \end{column}
        \begin{column}{0.48\textwidth}
            \textbf{Significance:}
            \begin{itemize}
                \item Standardized dataset for comparison
                \item Drives progress in medical AI
                \item Foundation for clinical translation
            \end{itemize}

            \vspace{0.2cm}
            \begin{block}{\footnotesize Relevance for Team Project}
                \footnotesize
                BraTS is directly relevant to the glioblastoma project - students can reference BraTS architectures.
            \end{block}
        \end{column}
    \end{columns}
\end{frame}

% =====================================================================
% SECTION: Quantitative Imaging
% =====================================================================
\section{Quantitative Imaging}

\begin{frame}{B05: Describe radiomic features and quantitative imaging}
    \textbf{\href{https://pyradiomics.readthedocs.io/}{Radiomics} = quantification of image characteristics}
    
    \vspace{0.3cm}
    \textbf{What are radiomic features?}
    \begin{itemize}
        \item Mathematical description of texture, shape, intensity in images
        \item Hundreds of features from a single region of interest (ROI)
        \item ``Deep'' properties not visible to the human eye
    \end{itemize}
    
    \vspace{0.3cm}
    \textbf{Feature categories:}
    \begin{columns}[T]
        \begin{column}{0.32\textwidth}
            \textbf{Intensity:}
            \begin{itemize}
                \item Mean
                \item Standard deviation
                \item Histogram features
            \end{itemize}
        \end{column}
        \begin{column}{0.32\textwidth}
            \textbf{Texture (GLCM):}
            \begin{itemize}
                \item Homogeneity
                \item Contrast
                \item Entropy
            \end{itemize}
        \end{column}
        \begin{column}{0.32\textwidth}
            \textbf{Shape:}
            \begin{itemize}
                \item Volume
                \item Surface area
                \item Sphericity
            \end{itemize}
        \end{column}
    \end{columns}
    
    \vspace{0.3cm}
    \begin{block}{Clinical application}
        Radiomics can predict molecular markers (IDH status), prognosis, and treatment response - based solely on standard images.
    \end{block}
\end{frame}

% =====================================================================
% SECTION: Tools
% =====================================================================
\section{Tools for Medical Image Analysis}

\begin{frame}{B06: Know about nnU-Net and MONAI as tools}
    \begin{columns}[T]
        \begin{column}{0.48\textwidth}
            \textbf{nnU-Net:}
            \begin{itemize}
                \item ``No-new-Net'' -- self-configuring
                \item Analyzes data and selects optimal configuration
                \item Preprocessing, architecture, training
                \item State-of-the-art on many benchmarks
                \item Only requires data -- the rest is automatic
            \end{itemize}
            
            \vspace{0.3cm}
            \textbf{Strengths:}
            \begin{itemize}
                \item Minimal manual tuning
                \item Reproducible baseline
            \end{itemize}
        \end{column}
        \begin{column}{0.48\textwidth}
            \textbf{\href{https://monai.io/}{MONAI}:}
            \begin{itemize}
                \item Medical Open Network for AI
                \item PyTorch-based framework
                \item Specialized for medical imaging
                \item Transformations, networks, metrics
                \item Active developer community
            \end{itemize}
            
            \vspace{0.3cm}
            \textbf{Strengths:}
            \begin{itemize}
                \item Flexibility for research
                \item Pre-built components
            \end{itemize}
        \end{column}
    \end{columns}
    
    \vspace{0.2cm}
    \begin{block}{\footnotesize Other tools}
        \footnotesize \textbf{3D Slicer:} Visualization and analysis | \textbf{ITK/SimpleITK:} Image processing | \textbf{Nipype:} Neuro-pipelines
    \end{block}
\end{frame}

% =====================================================================
% SUMMARY
% =====================================================================
\section*{Summary}

\begin{frame}{Summary: Medical Image Analysis}
    \textbf{Key points:}
    \begin{itemize}
        \item \textbf{B01:} MRI principles -- magnetization, RF pulse, relaxation, detection
        \item \textbf{B02:} MR sequences -- T1, T1+Gd, T2, FLAIR, DWI
        \item \textbf{B03:} Segmentation -- classify each pixel, U-Net, nnU-Net
        \item \textbf{B04:} BraTS -- benchmark for brain tumor segmentation
        \item \textbf{B05:} Radiomics -- quantitative features for prediction
        \item \textbf{B06:} Tools -- nnU-Net (automatic), MONAI (flexible)
    \end{itemize}

    \vspace{0.2cm}
    \begin{block}{\footnotesize Team project relevance}
        \footnotesize Medical image analysis is at the core of the glioblastoma team project. Apply your knowledge of MRI, segmentation, and radiomics.
    \end{block}

    \begin{block}{\footnotesize Lab 2}
        \footnotesize Hands-on experience with CNNs for image analysis, including medical applications (MR dementia).
    \end{block}
\end{frame}

\end{document}
